% Copyright (c) 2026 Eugene Joseph M. Ragasa.
% SPDX-License-Identifier: Apache-2.0
% Author: Eugene Joseph M. Ragasa
% Formal proof completed: 2026-10-31
\documentclass[11pt]{article}
\usepackage[T1]{fontenc}\usepackage{amsmath,amssymb}\usepackage[margin=1in]{geometry}\usepackage[hidelinks]{hyperref}
\title{Proof Dossier: Pointwise-Unitary Triple Convolution}\author{Eugene Joseph M. Ragasa}\date{Formal proof completed: 2026-10-31}
\begin{document}\maketitle
\section{Identity and status}
ID \texttt{BG-FOURIER-GAUGE-THM-001}; status \texttt{proved}. Lean declaration:
\par\noindent{\tiny\path{ProjectKoios.BandGauge.Fourier.Gauge.}}\\
{\tiny\path{finiteFourierCoefficient_pointwiseUnitaryConjugation_convolution}}.
\section{Exact statement and proof}
For a finite abelian momentum group, a unitary matrix field $U$, and a square matrix field $A$,
\[
 \widehat{U^\dagger A U}
 =\bigl((p\mapsto \widehat U(-p)^\dagger)\star\widehat A\bigr)\star\widehat U.
\]
Equivalently, the coefficient at $m$ is a finite sum of products
$\widehat U(a)^\dagger\widehat A(b)\widehat U(c)$ whose indices satisfy
$-a+b+c=m$. Conjugate transpose negates the Fourier mode under the repository convention, and two
applications of the finite matrix-product convolution theorem give the result.
\section{Boundary}
Unitarity specifies the existing pointwise gauge carrier, but the coefficient multiplication law is
algebraic. The theorem does not select or authenticate a gauge, bind runtime coefficients, establish
optimizer behavior, infer continuum decay, or assign a physical hopping range.
\section{Trusted base and verification}
Checked with Lean \texttt{4.35.0} and Mathlib revision
\texttt{56a005855ea17e0b0958f73ec801ae140b0a1679}. Kernel axioms are
\texttt{propext}, \texttt{Classical.choice}, and \texttt{Quot.sound}; no admissions.
AI-assisted and kernel-checked; no independent expert review recorded.
\end{document}
